\chapter{Carr and Kennedy (1994)}
\label{chap:ck94}

% Paper-specific labels: chap:/sec:/eq:/lst:/fig:ck94-<topic>.
\section{Array References}
\label{sec:ck94-array-references}

We distinguish between a \textbf{static} array reference and
an \textbf{execution} of an array reference.

A \textbf{static reference} maps an iteration vector (a point in the
iteration space) to an array element. With $n$ enclosing loops and $r$
array dimensions, we write
\begin{equation}
  u: \mathbf{I} \in \mathbb{Z}^{n} \longmapsto
  A[f_{1}(\mathbf{I})]\cdots[f_{r}(\mathbf{I})].
  \label{eq:ck94-static-reference}
\end{equation}
For a two-dimensional array, this becomes
$(i,j) \mapsto A[f_{1}(i,j)][f_{2}(i,j)]$.

Let $\mathbb{M}$ denote the set of static memory references, so
$u \in \mathbb{M}$.

It can be expressed in a more compressed form as a \textbf{matrix
transformation}.
\todo{define matrix transformation}

A \textbf{dynamic execution} of a reference is obtained giving the concrete value of the
index vector.
For example, when $i=1, j=2$, the dynamic execution of $u \coloneqq A[i][j]$ is
$A[1][2]$, mathematically, it is denoted as $u@(1,2)$.

\subsection{Variant and Invariant References}

A \textbf{single induction variable (SIV) subscript} reference contains
at most one loop induction variable in each subscript position.

\textbf{Variant} and \textbf{invariant} memory refereces are subtypes of SIV
subscript memory references. They are the most prominent types of references
\cite{carrKennedy1994}.

Their dependencies can be easily constructed and their distance vectors can
be found using \cref{eq:ck94-unrolled-distance}, which can be used as part of
the objective function in optimisaing loop balance.

A loop nest is denoted as $\Lambda=(\lambda_1,\ldots,\lambda_N)$, where
$\lambda_r$ is the loop at nesting level $1 \le r \le N$, numbered
from outermost ($\lambda_1$) to innermost ($\lambda_N$). When its
induction-variable names are distinct, let $\operatorname{IV}(\Lambda)$ be
their set and define
$\ell_{\Lambda}:\operatorname{IV}(\Lambda)\to\{1,\ldots,N\}$ by
$\ell_{\Lambda}(x)=r$ if $x$ is the induction variable of $\lambda_r$.
The mapping is local to $\Lambda$: the same name can denote a different level
in another nest. Thus a loop named by its induction variable $x$ is
$\lambda_{\ell_{\Lambda}(x)}$; subscripts on $\lambda$ always denote nesting levels.

\begin{lstlisting}[label={lst:ck94-variant-invariant-references},
language=C,caption={Loop Variant and Invariant References.}]
for (int i = 0; i < N; i++) {
  for (int j = 0; j < N; j++) {
      A[i][j] = B[i] * C[j];
  }
}
\end{lstlisting}


In \cref{lst:ck94-variant-invariant-references},
$\ell_{\Lambda}(i)=1$ and $\ell_{\Lambda}(j)=2$. Thus $B[i]$ is invariant
with respect to $\lambda_{\ell_{\Lambda}(j)}=\lambda_2$, while $C[j]$ is
invariant with respect to $\lambda_{\ell_{\Lambda}(i)}=\lambda_1$.

If a specific reference $v$ is invariant w.r.t. a loop $\lambda_{i}$, then
$\lambda_{i}$ is $v$-invariant.

If there exsists any reference that is invariant w.r.t a loop $\lambda_{i}$,
$\lambda_{i}$ is reference-invariant.

So $\lambda_1$ is $C[j]$-invariant and $\lambda_2$ is
$B[i]$-invariant. Both $\lambda_1$ and $\lambda_2$ are reference-invariant.

\section{Reference Dependency Representation}
Here we concern only about two memory references that accesses the
\textbf{same memory location}.

For \textbf{two references $u$ and $v$ contained in $n$ common loops}, seperate
executions of these references can be described with an \textit{iteration
vector}
$\mathbf{I} \in \mathbb{Z}^{n}$.

\subsection{Distance Vector}
Suppose executions $u@(\mathbf{I}_{u})$ and $v@(\mathbf{I}_{v})$
access the same element. Their \textbf{distance vector} (from $u$ to $v$) is
\begin{equation}
  \mathbf{d}_{u\to v} = \mathbf{I}_{v} - \mathbf{I}_{u}.
  \label{eq:ck94-distance-vector}
\end{equation}


Here, $u$ and $v$ must reference the same array, we do \textbf{not}
consider references to two different arrays because they do not reult benefits
in cache hit.

Note that even for two references with deterministic iteration index to memory
address functions, the distance vector may not be unique.

In \cref{lst:ck94-nonunique-distance}, we have $u \coloneqq A[2 \cdot i]$ and $v \coloneqq
A[i]$. The distance vectors of $u \rightarrow v$ dependency forms a set
% This is an instance of a definition for one example, so leave it unnumbered.
\[
\mathcal{D}_{u \rightarrow v}
= \{(i): 0 \le i < N,\; 2i < N\}.
\]
Here the write $u$ occurs at iteration $i$ and the matching read $v$ at
iteration $2i$; both iterations must be within the loop bounds.

\begin{lstlisting}[label={lst:ck94-nonunique-distance},
language=C, caption={Non-unique Distance Vector Example}]
for (int i = 0; i < N; i++) {
  A[i * 2] /* u */  = /* ... */;
  /* ... */         = A[i] /* v */;
}
\end{lstlisting}

In general, for any dependency execution $u \rightarrow v$, the set of
dependency vector is defined as
\begin{equation}
\mathcal{D}_{u \rightarrow v}
  = \{\mathbf{I}_{v} - \mathbf{I}_{u}:
  u@(\mathbf{I}_{u}) = v@(\mathbf{I}_{v})\}.
  \label{eq:ck94-distance-vector-set}
\end{equation}

By \cref{eq:ck94-distance-vector}, the dependence vectors between two references
can be calculated either way.
\[
\mathbf{d}_{u \rightarrow v}= - \mathbf{d}_{v \rightarrow u}
\]

Hence we cannot be sure about that \textit{without a iteration order}
for two references, e.g. a write reference $W$ and a read reference $R$,
whether this is a WAR or RAW dependency.

In a \textbf{normalised sequential loop nest}, iterations execute in
lexicographic order. With the destination-minus-origin convention of
\cref{eq:ck94-distance-vector}, this gives
\begin{equation}
  \begin{cases}
    u \rightarrow v &\text{if } \mathbf{d}_{u\to v} \succ_{\mathrm{lex}} \mathbf{0}, \\
    v \rightarrow u &\text{if } \mathbf{d}_{u\to v} \prec_{\mathrm{lex}} \mathbf{0}, \\
    \text{inspect statement order} &\text{if } \mathbf{d}_{u\to v} = \mathbf{0}.
  \end{cases}
  \label{eq:ck94-dependence-order}
\end{equation}
The first nonzero component of $(0,0,-1,1)$ is negative, so the
inter-iteration ordering is $R\rightarrow W$ if this vector was computed
as $\mathbf{d}_{W\to R}$. The access types alone do not establish RAW.

When $\left(\mathbf{d}_{u \rightarrow v}\right)_{i} = 0$, we say that the edge
$e = \left\langle u, v \right\rangle$ is invariant w.r.t. loop $\lambda_{i}$,
denoted as
\[
  Invar(e, \lambda_{i})
\]

\subsection{Direction Vector}
The direction vector $\mathbf{D}$ can be a more compressed (but lossy) representation of the
distance vector set $\mathcal{D}$.
\begin{equation}
  \label{eq:ck94-direction-vector}
\mathbf{D}_{i}=
\begin{dcases}
  \ge & \text{if } \forall \mathbf{d} \in \mathcal{D}: d_{i} \ge 0 \\
  > & \text{if } \forall \mathbf{d} \in \mathcal{D}: d_{i} > 0 \\
  = & \text{if } \forall \mathbf{d} \in \mathcal{D}: d_{i} = 0 \\
  < & \text{if } \forall \mathbf{d} \in \mathcal{D}: d_{i} < 0 \\
  \le & \text{if } \forall \mathbf{d} \in \mathcal{D}: d_{i} \le 0 \\
  \neq & \text{if } \forall \mathbf{d} \in \mathcal{D}: d_{i} \neq 0 \\
  * & \text{otherwise}
\end{dcases}
=
\begin{dcases}
  \ge & \text{if } \forall \mathbf{I}_{u}^{*}, \mathbf{I}_{v}^{*}:
  \mathbf{I_{v}}^{*}_{i} \ge \mathbf{I_{u}}^{*}_{i} \\
  > & \text{if } \forall \mathbf{I}_{u}^{*}, \mathbf{I}_{v}^{*}:
  \mathbf{I_{v}}^{*}_{i} > \mathbf{I_{u}}^{*}_{i} \\
  = & \text{if } \forall \mathbf{I}_{u}^{*}, \mathbf{I}_{v}^{*}:
  \mathbf{I_{v}}^{*}_{i} = \mathbf{I_{u}}^{*}_{i} \\
  < & \text{if } \forall \mathbf{I}_{u}^{*}, \mathbf{I}_{v}^{*}:
  \mathbf{I_{v}}^{*}_{i} < \mathbf{I_{u}}^{*}_{i} \\
  \le & \text{if } \forall \mathbf{I}_{u}^{*}, \mathbf{I}_{v}^{*}:
  \mathbf{I_{v}}^{*}_{i} \le \mathbf{I_{u}}^{*}_{i} \\
  \neq & \text{if } \forall \mathbf{I}_{u}^{*}, \mathbf{I}_{v}^{*}:
  \mathbf{I_{v}}^{*}_{i} \neq \mathbf{I_{u}}^{*}_{i} \\
  * & \text{otherwise}
\end{dcases}
\end{equation}
So we can have $\mathbf{D}_{u \rightarrow v} = \langle \ge \rangle$.

\subsection{Carrier Loop}
The \textbf{\textit{carrier loop} of this dependency instance} is the loop
corresponds to the \textbf{first non-zero component} of the distance vector
$\mathbf{d}$.

For example, assume the loop nest has induction variables
$\langle i,j,k\rangle$ at levels $\langle 1,2,3\rangle$.
If $u@\mathbf{I_{u}}_{p} = v@\mathbf{I_{v}}_{q}$, and
\[
\mathbf{I_{v}}_{q} - \mathbf{I_{u}}_{p} = <0,
\underbrace{-1}_{\text{first }\ne 0}, 1>
\]
then loop $\lambda_2$ (whose induction variable is $j$)
is the carrier loop of this specific instace of the dependency $u \rightarrow v$.
There is nothing preventing us to construct the dependency where
$\mathcal{D}_{u \rightarrow v} = \{(0,1,0),(1,0,0)\}$ where both loops
(i.e. $\lambda_1$ and $\lambda_2$) are the carrier loops for some instances
of $u \rightarrow v$.

The \textbf{threashold} for a carried dependence is the value of the first non-zero entry
of that distance vector. The \textbf{minimum threshold} is the \textbf{minimum
possitive} threashold.

If for a dependency, the threashold is constant throughout execution, this
dependency has \textbf{constant threashold}, and such dependency is a
\textbf{constant dependency}. In
\cref{lst:ck94-nonunique-distance}, the threshold is not constant; however,
for references $A[i]$ and $A[i+c]$, the threshold is constant (equals to $c$).
In \cite{carrKennedy1994}, only constant dependencies are considered.


% We need to distinguish between thhe oridnaly \textbf{computation/dataflow
% dependency} and \textbf{memory dependency between array references}, since in
% the conrext of data cache optimisation, we care only about cache hit, which has
% the same or adjacent memeory location.

% \begin{lstlisting}[language=C,caption={Memory references in a loop nest.},
%   label={lst:mem-dep-ex-1}]
%   for (int i = 1; i < N; i++) {
%     for (int j = 0; j < N; j++) {
%         A[i][j] /* ref W_A */ = 1 + B[i-1][j]; /* ref R_B  */
%         C[i][j] /* ref W_C */ = 2 * A[i][j]  ; /* ref R_A1 */
%         D[i][j] /* ref W_D */ =     A[i][j-1]; /* ref R_A2 */
%     }
%   }
% \end{lstlisting}
%
% In \Cref{lst:mem-dep-ex-1}, although we have the following dataflow dependencies,
% they do not reference the same memory location.
% \[
% R_{B} \xrightarrow[\text{add}]{\text{dataflow}} W_A \quad
% R_{A_{1}} \xrightarrow[\text{multiply}]{\text{dataflow}} W_{C} \quad
% R_{A_{2}} \xrightarrow[\text{assign}]{\text{dataflow}} W_{D}
% \]
%
% It is important to note that $R_{B} \xrightarrow{\text{dataflow}}
% W_A$ means "at iteration $(i,j)$, memory location $R_B@(i-1,j)$ is read". It
% \textbf{does not} mean that "iteration $(i,j)$ read the value produced by
% iteration", since $B$ is never produced within this loop net.
%
% On the contrary, in the same iteration, we have the following memory
% dependencies to the same memory location.
% \[
% \begin{aligned}
% &W_A \xrightarrow{\mathbf{d}(0,0)}R_{A_{1}} \\
% &W_A \xrightarrow{\mathbf{d}(0,1)}R_{A_{2}} \\
% &R_{A_{1}} \xrightarrow{\mathbf{d}(0,1)}R_{A_{2}}
% \end{aligned}
% \]

% the \textbf{same memory address}.
% \begin{lstlisting}[language=C]
% for (int i = 1; i < M; i++) {
%   for (int j = 1; i < N - 2; j++) { /* assume boundaries are handled correctly */
%     a[i][j]     /* W1 */ = A[i-1][j+1]; /* A1 */
%     A[i-1][j+2] /* W2 */ = A[i][j-1]  ; /* A2 */
%   }
% }
% \end{lstlisting}
%
% For example, in \todo{listing}, we solve for
% iterations $(i, j)$ that accesses the same location $A[p][q]$.
% \begin{table}[htbp]
% \centering
% \begin{tabular}{lcc}
% \hline
% \textbf{Reference} & \textbf{Address equation} & \textbf{Execution accessing $A[p][q]$} \\
% \hline
% \texttt{W1} & $(i, j) = (p, q)$ & $(p, q)$ \\
% \texttt{R1} & $(i - 1, j + 1) = (p, q)$ & $(p + 1, q - 1)$ \\
% \texttt{W2} & $(i - 1, j + 2) = (p, q)$ & $(p + 1, q - 2)$ \\
% \texttt{R2} & $(i, j - 1) = (p, q)$ & $(p, q + 1)$ \\
% \hline
% \end{tabular}
% \caption{Reference addresses and corresponding execution access equations.}
% \label{tab:reference_equations}
% \end{table}
%
% In \Cref{tab:reference_equations}
%
% The distance vector is defined as the difference between the read index (sink) and
% the write index (source).
% \[
% \begin{aligned}
% &\mathbf{i_{\text{sink}}} = \mathbf{i_{\text{source}}} + \mathbf{d} \\
% &\mathbf{i_{\text{read}}} = \mathbf{i_{\text{write}}} + \mathbf{d}
% \end{aligned}
% \]
%
% For references $R_{A_{2}}$ and $W_{A}$
%
% % Note that in iteration $(\hat{i}, \hat{j}-1)$, we write to location
% % $A[\hat{i}][\hat{j}-1]$ by reference $W_A \coloneqq A[i][j]$, and in
% % iteration $(\hat{i}, \hat{j})$, we read from \textbf{the same memory location}
% % $A[\hat{i}][\hat{j}-1]$ by refenrence
% % $R_{A_{2}} \coloneqq A[i][j-1]$. Hence we have the dependece.
% % $W_A \rightarrow R_{A_{2}}$
%
%
%
%
% \todo{Review this}
% We still keep the read-after-read dependence $R_{A_{1}}
% \xrightarrow{\mathbf{d}(0,1)}R_{A_{2}}$, treating $R_{A_{1}}$ as the source and
% $R_{A_{2}}$ as the sink. \todo{Why do Carr and Kennedy keep this?}
%
% Note that the distance victor for $W_A \xrightarrow{\mathbf{d}(0,0)}R_{A_{1}}$
% is $\mathbf{0}$, and we call this dependence \textbf{loop-independent}
% dependency, while the previous two are called loop-dependent dependencies.
%
% For loop-ependent dependencies, the \textbf{carrier loop} is the
% \textbf{first non-zero component} of its distancevector $\mathbf{d}$. \todo{add
% intuition and examples}

\section{Dependence Analysis}
\label{sec:ck94-dependence-analysis}
\Cref{lst:ck94-unroll-original,lst:ck94-unroll-expanded,lst:ck94-unroll-jammed}
illustrate unroll-and-jam (\cref{sec:ck94-unroll-and-jam}): expand an outer
loop, then fuse the resulting inner loops. We use this example to introduce
dependence analysis.
\begin{figure}[H]
  \centering
  \begin{minipage}[t]{0.30\textwidth}
\begin{lstlisting}[language=C, basicstyle=\ttfamily\scriptsize,
  caption={Original},
  label={lst:ck94-unroll-original}]
for (int i = 0; i < N; i++) {
  for (int j = 0; j < N; j++) {
    A[i][j] /* W0 */ = A[i-2][j] /* R0 */;
  }
}
\end{lstlisting}
  \end{minipage}%
  \hfill
  \begin{minipage}[t]{0.35\textwidth}
\begin{lstlisting}[language=C, basicstyle=\ttfamily\scriptsize,
  caption={Unrolled},
  label={lst:ck94-unroll-expanded}]
for (int i = 0; i < N; i += 3) { /* k=3 */
  for (int j = 0; j < N; j++) {
    A[i  ][j] /* W1 */ = A[i-2][j] /* R1 */;
  }
  for (int j = 0; j < N; j++) {
    A[i+1][j] /* W2 */ = A[i-1][j] /* R2 */;
  }
  for (int j = 0; j < N; j++) {
    A[i+2][j] /* W3 */ = A[i  ][j] /* R3 */;
  }
}
\end{lstlisting}
  \end{minipage}%
  \hfill
  \begin{minipage}[t]{0.31\textwidth}
\begin{lstlisting}[language=C, basicstyle=\ttfamily\scriptsize,
  caption={Unroll-and-Jam ($k = 3$)},
  label={lst:ck94-unroll-jammed}]
for (int i = 0; i < N; i += 3) { /* k=3 */
  for (int j = 0; j < N; j++) {
    A[i  ][j] /* W1 */ = A[i-2][j] /* R1 */;
    A[i+1][j] /* W2 */ = A[i-1][j] /* R2 */;
    A[i+2][j] /* W3 */ = A[i  ][j] /* R3 */;
  }
}
\end{lstlisting}
  \end{minipage}
\end{figure}

\subsection{General Case}
\Cref{lst:ck94-unroll-jammed} shows the transformation applied to
\cref{lst:ck94-unroll-original}. The iteration indices of the outer
loop $\lambda_1$ (induction variable $i$) and inner loop $\lambda_2$
(induction variable $j$) can be expressed as
% Coordinates and calculations for this worked example stay unnumbered.
\[
\begin{aligned}
&I_{1} = 3x &\quad \forall 0 \le x < \left\lceil \frac{N}{3} \right\rceil \\
&I_{2} = y &\quad \forall 0 \le y < N
\end{aligned}
\]
We collect the read set $R=\{R_{1},R_{2},R_{3}\}$ and write set
$W=\{W_{1},W_{2},W_{3}\}$. Boundary iterations require a remainder loop
or guards if $N$ is not divisible by three.

To analyse the read-write dependency, for $\langle R_{i}, W_{j} \rangle$ pair,
we solve the equation where they point to the same memory address. For example,
for $\langle R_{1}, W_{2} \rangle$, we solve for equations
\[
\begin{dcases}
  3 x_{R_{1}} - 2 &= 3 x_{W_{2}} + 1 \\
  y_{R_{1}} &= y_{W_{2}}
\end{dcases}
\implies
\begin{dcases}
  x_{R_{1}} - x_{W_{2}} &= 1 \\
  y_{R_{1}} - y_{W_{2}} &= 0
\end{dcases}
\implies
\mathbf{d}_{W_{2} \rightarrow R_{1}} = \langle 1, 0 \rangle.
\]
Thus $W_{2}$ writes the location read by $R_{1}$ in a later outer
iteration (provided both iterations are in range).

However, when we try to solve the dependency $\langle W_{1}, R_{2} \rangle$, we
have
\[
\begin{dcases}
  3 x_{W_{1}} &= 3 x_{R_{2}} - 1 \\
  y_{W_{1}} &= y_{R_{2}}
\end{dcases}
\implies
\begin{dcases}
  x_{W_{1}} - x_{R_{2}} &= \frac{1}{3} \\
  y_{W_{1}} - y_{R_{2}} &= 0
\end{dcases}
\implies
\mathbf{d}_{W_{1} \rightarrow R_{2}} = \left\langle \frac{1}{3}, 0 \right\rangle
\]
Which means we cannot find a valid (integral) iteration distance in which
references $W_{1}$ and $R_{2}$ access the same memory address. So there's no
dependency between $W_{1}$ and $R_{2}$.

Here, we defined the dependecy vector $\mathbf{d}_{W_{1} \rightarrow R_{2}} \in
\mathbb{Z}$. however, if we defined it as the distance between actual
iterations, it would be in $3\mathbb{Z}$ because the outer loop advances by three.
Either coordinate system is valid if used consistently; the fractional
solution above is not a valid pair of iterations.

\subsection{Unroll-and-Jam: Single Induction Variable Subscript References}
\label{sec:ck94-siv-dependences}

In this section, we discuss special dependencies \textbf{generated by
unroll-and-jam} from references with SIV subscripts.

If the $i$-th loop in a loop nest $\Lambda$ is \textbf{unrolled into $X_i$ copies}, the updated dependence graph
$G' = (V', E')$ can be computed from the original graph $G = (V,E)$.

Each vertex has $X_i$ copies, indexed from $0$ to $X_i-1$.
\[
\forall v \in V: \left\{v_{0}, \dots, v_{X_i-1}\right\} \subset V'.
\]

The dependencies (edges) are handled differently depending on whether
\textbf{the sink is loop variant or not}.

\subsubsection{Variant References}
For an edge $e$, we denote its distance vector as $\mathbf{d}(e)$, so its $i$-th
component (corresponding to the $i$-th loop) is $\mathbf{d}_{i}(e)$. Then
\[
  \forall e = \left\langle v, w \right\rangle \in E: \exists \left\{e_{0}, \dots,
  e_{X_i-1}\right\} \subset E^{\prime}
\]
where
\[
e_{m} = \left\langle v_{m}, w_{f(m;e,X_{i})} \right\rangle
\]
\begin{equation}
  f(m;e,X_{i}) = \left(m + \mathbf{d}_{i}(e)\right) \bmod X_{i}
\end{equation}
where $X_{i}$ is the number of loop bodies of loop $\lambda_{i}$ after
unrolling.

This means, that for every reference source $v_{i}$, there is always a
dependence edge to a sink at $w_{f(i)}$, and its distance vector
$\mathbf{d}(e_{i})$ is computed as

\begin{equation}
  \label{eq:ck94-unrolled-distance}
  \mathbf{d}_{i}(e_{m}) =
  \begin{dcases}
    \left\lfloor \frac{\mathbf{d}_{i}(e)}{X_{i}}\right\rfloor + 1
    &\text{ if } f(m; e, X_{i}) < m \\
    \left\lfloor \frac{\mathbf{d}_{i}(e)}{X_{i}}\right\rfloor
    &\text{ if } f(m; e, X_{i}) \ge m
  \end{dcases}
\end{equation}

This is a modulo $X_{i}$ system. A reference the $k$-th iteration of the
orignal loop will be mapped to the $(k \bmod X_{i})$-th copy in the
$\left\lfloor \frac{k}{X_{i}} \right\rfloor$-th iteration of the new loop.

\todo{why +1?}

\Cref{eq:ck94-unrolled-distance} shows that for the new dependency edge
$e_{m}=\left\langle v_{m}, w_{f(m)} \right\rangle$
to be loop-invariant regarding to $\lambda_{i}$, the original distace of
dependency edge $e = \left\langle v, w \right\rangle$ should be less than the
repetition after expansion. Otherwise, $\mathbf{d}_{i}(e) \ge X_{i}$
for a source $v_{m}$, its matching sink
$w_{f(m)}$ occours in later iterations instead of the same iteration.

\begin{equation}
    \mathbf{d}_{i}(e_{m}) = 0
    % \xLeftrightarrow{\text{only in case 2}} \;
    \iff
    \left(f(m; e, X_{i}) \ge m\right) \land \left(\left\lfloor
    \frac{\mathbf{d}_{i}(e)}{X_{i}} \right\rfloor = 0 \right)
  \label{eq:ck94-zero-distance-criterion}
\end{equation}

\begin{equation}
  \left\lfloor  \frac{\mathbf{d}_{i}(e)}{X_{i}} \right\rfloor = 0
  \iff 0 \le \frac{\mathbf{d}_{i}(e)}{X_{i}} < 1
  \xLeftrightarrow{X_{i} > 0} \;
  0\le \mathbf{d}_{i}(e) < X_{i}
  \label{eq:ck94-distance-below-factor}
\end{equation}

From \cref{eq:ck94-zero-distance-criterion,eq:ck94-distance-below-factor}, if
given $0 \le \mathbf{d}_{i}(e) < X_{i}$, we have
\begin{equation}
  \mathbf{d}_{i}(e_{m}) = 0 \iff f(m; e, X_{i}) \ge m
    \label{eq:ck94-zero-distance-no-wrap}
\end{equation}

Given $0 \le \mathbf{d}_{i}(e) < X_{i}$, we can prove that
\footnote{\todo{prove later}}
\begin{equation}
  f(m; e, X_{i}) \ge m \iff m < X_{i} - d_{i}(e)
  \label{eq:ck94-no-wrap-bound}
\end{equation}

With \cref{eq:ck94-no-wrap-bound,eq:ck94-zero-distance-no-wrap},
we have that: Given $0 \le d_{i}(e) < X_{i}$,
\begin{equation}
  \mathbf{d}_{i}(e_{m}) = 0 \iff m < X_{i} - \mathbf{d}_{i}(e)
  \label{eq:ck94-zero-distance-copy-bound}
\end{equation}

\Cref{eq:ck94-zero-distance-copy-bound} says that, for an original edge $e$, its
distance w.r.t. loop $\lambda_{i}$ is less than the loop's unroll factor, then
as long as $m$-th edge copy in the unroll-and-jam'ed loop has $m < X_{i} -
\mathbf{d}_{i}(e)$, its corresponding distance component is 0. This will be
useful when we are estimating how many references can be eliminated using scalar
replacement.


% Without loss of generality, we assume the following loop nest
% \begin{lstlisting}[label={lst:example},
% language=C,caption={Code example.}]
% for (int i = 0; i < M; i += k) { /* const int k */
%   for (int j = 0; j < N; j++) {
%     /* ----- Original body ------ */
%     /* loop body 0   */
%     /* --- Expanded bodies, k --- */
%     /* loop body 1   */
%     /* loop body ... */
%     /* loop body k   */
%   }
% }
%
% \end{lstlisting}
%
% \todo{say something about the old vs new $i$ index}
% \begin{equation}
%   \begin{aligned}
%     \mathbf{I} &= \left\langle x, y \right\rangle \quad
%     \left(0 \le x < M\right) \\
%     \mathbf{I}^{\prime} &= \left\langle kx^{\prime}, y^{\prime} \right\rangle
%     \quad \left(0 \le x^{\prime} < \left\lfloor \frac{M}{k} \right\rfloor \right)
%   \end{aligned}
% \end{equation}
%
% In a new iteration $\mathbf{I}^{\prime}$, the $m$-th copy of reference $v$
% ($v_{m})$ comes form the original iteration
% \[
%   (k+1)x + m
% \]

\subsubsection{Invariant References}
When applying unroll-and-jam, invariant references behave differently. Note that
here when we say a reference is \textit{variant}, it means variant w.r.t every loop,
otherwise it is considered invariant.

For an invariant reference w.r.t. the $\lambda_{k}$ loop, by definition, varying its
iteration coordinate $I_k$ of loop $\lambda_k$ does not change the memory location. Thus
the $k$-th component of its direction vector $\mathcal{D}_{k} =$ '$\neq$'.

After unrolling, we would have multiple instances of this $\lambda_k$-invariant reference, and
hence the diatance vector between the these references may result in $D_{l}=$ '='.

\textit{This would be problematic when it's in the carrying loop's position.}
\begin{figure}[H]
  \centering
  \begin{minipage}[t]{0.48\textwidth}
\begin{lstlisting}[language=C,
  caption={Original},
  label={lst:ck94-invariant-original}]
for (int i = 0; i < N; i++) {
  for (int j = 0; j < N; j++) {
    for (int k = 0; k < N; k++) {
      /* ... */ = A[i][k] /* R0 */
    }
  }
}
\end{lstlisting}
  \end{minipage}%
  \hfill
  \begin{minipage}[t]{0.48\textwidth}
\begin{lstlisting}[language=C,
  caption={Unrolled},
  label={lst:ck94-invariant-unrolled}]
for (int i = 0; i < N; i++) {
  for (int j = 0; j < N; j+=2) { /* <-- Unrolled by k = 2 */
    for (int k = 0; k < N; k++) {
      /* ... */ = A[i][k] /* R1 */;
      /* ... */ = A[i][k] /* R2 */;
    }
  }
}
\end{lstlisting}
  \end{minipage}
  \caption{Invariant Consistent Reference Example}
  \label{fig:ck94-invariant-unroll}
\end{figure}

For example, reference $R_{0}$ in \cref{lst:ck94-invariant-original} is invariant
w.r.t. $\lambda_2$ (induction variable $j$). When analysing dependencies,
we can solve the equation
\[
  \begin{dcases}
    I_1 = I'_1 \\
    I_3 = I'_3
  \end{dcases}
  \implies
  \begin{dcases}
    d_1 = I'_1 - I_1 = 0 \\
    d_3 = I'_3 - I_3 = 0
  \end{dcases}
  \implies
  \mathbf{d} = \left\langle 0, I'_2 - I_2, 0
  \right\rangle
\]
By imposing lexigraphical (normalised sequential loop) order of the outerloop,
we have $I'_2 > I_2$ for a forward dependence (we do \textit{not}
allow $I'_2 = I_2$ because we need \textit{distinct} accesses).
so its corresponding direction vector
\[
  \mathbf{D}_{R_{0}\rightarrow R_{0}} = \left\langle =, >, = \right\rangle
\]

It becomes problematic when, after replication, the \textbf{outermost non-zero} direction vector entry
(assume $l$) becomes zero.\[
  \begin{aligned}
  &\exists \mathbf{d}\in \mathcal{D}_{R_{0}\rightarrow R_{0}}&&: \mathbf{d}_{l}
  \neq 0 \\
  \land \quad
  &\exists \mathbf{d}, \mathbf{d}^{\prime} \in \mathcal{D}_{R_{1}\rightarrow
  R_{2}}&&:
  \mathbf{d}_{l} = 0 \land \mathbf{d}^{\prime}_{l} \neq 0
  \end{aligned}
\]

This would \textbf{change the carrying loop} when $\mathbf{d}_{l} = 0$.

Since we want to reuse \cref{eq:ck94-unrolled-distance}, however, in the
variant case, we have $\left| \mathcal{D}_{\text{invar}} \right| = 1$, so it is obvious to

\todo{Complete this}

\subsection{Unroll-and-Jam: Multiple Induction Variable Subscript References}
% In \cite{carrPhd}, references with multiple induction variable (MIV) subscript is discussed.
% \todo{Add here?}

\cite{carrPhd} supports a subset of MIV subscript. \todo{perhaps explain?}

\section{Unroll-and-Jam}
\label{sec:ck94-unroll-and-jam}
This method unrolls the inner loop and jams the resulting inner loops back
together. In \cite{carrKennedy1994}, unroll-and-jam is only to create more reuses.

In this report, we assume loops are unrolled in an innermost to outermost order
\cite{carrKennedy1994}. We can recursively unroll at an arbitrary
position by unrolling all its nesting loops. This order makes sense since we are
optimising to cache hit.


\textit{To generalise the unrolling factor $k$, we use \textbf{$X_{i}$} to express the unroll
factor at the $i$-th outer loop. In simple examples where we only unroll one
loop, $k$ is used.}

\subsection{Safety}
Unroll-and-Jam is safe if the flow of value is not changed, i.e. it protects the
original ordering of dynamic operations
The safety theorem from \todo{cite callahan et al. 1998} states that unrolling
loop $k$ more than ...

Let edge $e$ be a true (RAW), output (WAW) or antidependence (WAR)
\textit{carried at level $k$} with dependence vector
\[
  \mathbf{d} = \left\langle 0, \dots, d_{k}, 0, \dots, 0, d_{j}, \dots
  \right\rangle
\]
where $d_{j} < 0$ and $\forall k < i < j: (d_{i} = 0)$, then the loop at level k
($\lambda_{k}$) can be unrolled \textbf{at most} $d_{k}-1$ times before a
dependence is generated that prevents fusion of the inner $n-k$ loops. Also,
\begin{itemize}[nosep]
  \item \textbf{Loop fusion} can be introduced before unroll-and-jam safty is
    checked to remove negative entires in preventing dependencies \todo{Wolfe
    1986};
  \item \textbf{Loop distribution} must be safe when non-do statements appear at
    levels non-innermost levels or when multiple loops appear at inner levels.
\end{itemize}



\section{Objective Function}
% The derivation assumes
% \footnote{For clarity of presentation, not the limitation of the algorithm
% \cite{carrKennedy1994}}:
% \begin{enumerate}[nosep]
%   \item Perfectly-nested loop nets \todo{section 3.2.4.5 \cite{carrKennedy1994}}
%   \item Loops with conditional control in the innermost loop body is discarded as
%     a unroll-and-jam candidate
%     \footnote{Inserting controlflow have disastrous effects upon performance
%     \cite{carrKennedy1994}}
%   \item Scalar replacement removes no stores
%   \item At most one incoming or outgoing edge incident upon each $v \in V$
%     \cite{carrKennedy1994}, whenever this is assumed, the function $e: V \mapsto
%     E$ is used, so the only-edge of reference $v$ is calculated as $e(v)$. For
%     simplicity, $e_{m}(v)$ is used to denote the $m$-th copy of the original
%     edge $e(v)$, also, $e_{0}(v) \coloneqq e(v)$.
% \end{enumerate}

\textbf{Machine Balance} is a compile-time constant defined by the target
architecture.
\begin{equation}
  \label{eq:ck94-machine-balance}
\beta_{M} = \frac{M_{M}}{F_{M}}
\end{equation}
Where $M_{M}$ is the rate of fetching data from memeory
% These units explain a model parameter; the model itself is numbered.
\[
M_{M} = \frac{\text{max words}}{\text{cycle}}
\]
and $F_{M}$ is the aate of floating-point operations can be performed
\[
F_{M} = \frac{\text{max flops}}{\text{cycle}}
\]
\textbf{Loop Balance} is determined at compile time for a specific loop
\begin{equation}
  \label{eq:ck94-loop-balance}
\beta_{L} = \frac{M}{F}
\end{equation}
where $M$ is the number of memory references and $F$ is the number of
floating-point operations.

\paragraph{Optimisation Objective}
The objective of this optimisation is to find two outer loops $\lambda_{j},
\lambda_{k}$ to unroll
\footnote{Unroll-and-jam can be limited to one or two loops in a given loop nest
\cite{carrKennedy1994}. Most loops are not very deep.}, and find their
corresponding unroll factors that
minimises loop balance while respecting the available register count.

\textbf{Objective function.} Minimise memory operations per floating-point
operation, subject to the register constraint:
% \begin{equation}
%   \min_{X_i,X_j}\;\beta_L(X_i,X_j)
%   = \min_{X_i,X_j}\;\frac{M(X_i,X_j)}{F(X_i,X_j)}.
% \end{equation}
\begin{equation}
  \min_{X_i}\;\beta_L(X_i)
  = \min_{X_i}\;\frac{M(X_i)}{F(X_i)}.
\end{equation}

\textbf{Register constraint.}
\begin{equation}
  % R(X_{i}, X_{j}) \le R_{\text{arch}}
  R(X_{i}) \le R_{\text{arch}}
  % \label{eq:}
\end{equation}

Unroll-and-jam together with scalar replacement can reduce memory access.
However, using scalar replacement will introduce intermediate registers, hence
the constraint must be met.

% \subsection{Simplification and Generalised Cases}
% \subsection{Multiple Edge References}
% We first consider a spetial case where every reference has only one incident
% (be it incoming or outgoing) edge \todo{exclusive or?},
% and we denote that edge related to that
% referece $v$ as $e(v)$, with $e: V \mapsto E$.
%
% In the formation of our loop balance objective function, distance vector is used.
% In the simplified case, we can associate exactly one distance vector to
% this reference
% \[
%   \mathbf{d}(v) \coloneqq \mathbf{d}(e(v))
% \]
% In general, references have different edges, but the idea is to find a
% \textbf{representative} distance vector for this reference so we can treat this
% reference
%
% \subsection{Imperfectly-nested Loop Nets}
% \todo{}

\section{Floating-point Operations}
The number of floating-point operations $F$ is calculated as
\begin{equation}
  \label{eq:ck94-flop-count}
  F(X_{i}) = f \cdot \prod_{1 \le i < N} X_{i}
\end{equation}
Where $f$ is the number of floating-point operations in the original loop,
$X_{i}$ is the number of loop bodies after unroll-and-jam (expansion factor
$+1$). This formula assumes that floating-point common subexpressions are not
removed from the inner loop \cite{carrKennedy1994}. \todo{Explain possible
situations where such common subexpressions are removed}


\section{Memory References}
\todo{Make reference partition a seperate section since it will be used in
register pressure modelling as well}
When calculating the number of memory references, \textbf{infinite register set}
is assumed \cite{carrKennedy1994}.

A specific reference set $V \subset \mathbb{M}$ can be partitioned into the
following subsets
\begin{itemize}
  \item $V^{\varnothing}$: \textit{references} without incoming consistent dependence;
  \item $\{V^{C}_r$, $V^{C}_w\}$\footnotemark[\value{footnote}]: 
    \textit{memory reads and writes} with  incoming 
    (loop-carried or loop-independent) consistent dependence 
    but is \textbf{variant w.r.t. some loop} (not invariant w.r.t. any loop);
  \item $\{V^{I}_{r}$, $V^{I}_w\}$\footnotemark[\value{footnote}]: 
    \textit{memory reads and writes} \textbf{invariant w.r.t some loop}.
\end{itemize}
\footnotetext{\cite{carrKennedy1994} only mentioend the sets $\{V^{\varnothing},
V^{C}_r, V^{I}_r\}$. \cite{carrPhd} added $\{V^{C}_w, V^{I}_w\}$}

\subsection{Non-dependent References}
For references without incoming consistent dependencies, unrolling always
introduces memeory access, hence the total memory accesses after unrolling is
\begin{equation}
  M^{\varnothing}(\mathbf{X}) = \sum_{v \in V^{\varnothing}}
  \left(\prod_{1 \le i < N} X_{i} \right)
\end{equation}

For a specific loop $\lambda_{i}$, unrolling it have memory accesses
\begin{equation}
  \begin{aligned}
    M^{\varnothing}(X_{i}; \lambda_{i}) &= \sum_{v \in V^{\varnothing}} X_{i} \\
    &= \left| V^{\varnothing} \right| X_{i} \\
    &= a_{0} X_{i}
  \end{aligned}
\end{equation}

\subsection{Variant References}
$V^{C}_{r}$ is of particular interest of \textbf{scalar replacement},
since it can reduce the number of memory accesses of references,
orignally having cross-iteration  dependeicies,
that happen to be in the same iteration after unroll-and-jam.

\subsubsection{Special Case: Single-edge References}
For scalar replacement to remove memory access for variant references, it must
be of form $\mathbf{d} = \left\langle 0, 0, 0, \dots, \mathbf{d}_{N} \right\rangle$. i.e. the
dependency edge must be carried by the innermost loop or loop-independent.

From \cref{eq:ck94-zero-distance-copy-bound}, since now we have the one-to-one
correspondance between reference $v$ and its only \textit{incoming} edge $e(v)$,
we can find all copies with $\mathbf{d}_{i}(e(v)) = 0$.
\[
  \mathcal{C}(v; X_{i}) \coloneqq
  \left\{ v_{m}: m < X_{i} - \mathbf{d}_{i}(e(v)) \right\}
\]

Essentially, the \textbf{number of edges} that are \textit{amenable to scalar replacement
is}
\begin{equation}
  \begin{aligned}
    \left| \mathcal{C}(v; X_{i}) \right|
    &= m \\
    &= \max \left(X_{i} - \mathbf{d}_{i}(e(v)), \; 0 \right) \\
    &\xlongequal{x^{+} \coloneqq \max(x, 0)}
    \left(X_{i} - \mathbf{d}_{i}(e(v))\right)^{+}
  \end{aligned}
  \label{eq:ck94-removable-copies}
\end{equation}



According to \cref{eq:ck94-removable-copies}, the total number of
memory references \textbf{removed by scalar replacement} is calculated as
% \todo{are we assuming single edge $e(v)$ to vertex $v$? Yes we do
% \cite{carrKennedy1994}}
\begin{equation}
  \sum_{v \in V^{C}_{r}} \left(\prod_{1 \le i < N}
    \left(X_{i} - \mathbf{d}_{i}(e(v))\right)^{+}\right)
  \label{eq:ck94-removed-variant-references}
\end{equation}

Note that \cref{eq:ck94-removed-variant-references} implies that,
for an edge \(e\), if \textit{any} outer loop $\lambda_{i}$ satisfies \(d_i(e)\ge X_i\),
then no copies of the sink reference can have their memory accesses eliminated
through \(e\) by the modeled scalar replacement under the chosen unroll factors.

\[
  \exists \lambda_{i}: X_i\le d_i(e)
\implies (X_i-d_i(e))^+=0
\implies \prod_{1 \le j < N}(X_j-d_j(e))^+=0.
\]

\textit{Since we need the two references to be \textbf{reachable in any dimensions}, as
long as in one dimension, it is not reachable, the two references are irrechable
through the distance vector $\mathbf{d}(e)$.}

The total number of memory references \textbf{before scalar replacement} is
calculated as
\begin{equation}
  \sum_{v \in V^{C}_{r}} \left(\prod_{1 \le i < N} X_{i}\right)
  \label{eq:ck94-variant-references-before-replacement}
\end{equation}

With  total number of memory references \textbf{after scalar replacement} is
calculated as
\begin{equation}
  M^{C}_{r}(\mathbf{X}) =
  \sum_{v \in V^{C}_{r}} \left(
  \prod_{1 \le i < N} X_{i} \; - \;
  \prod_{1 \le i < N} \left( X_{i} - \mathbf{d}_{i}(e(v))\right)^{+}
  \right)
  % \label{eq:}
\end{equation}

For a specific loop $\lambda_{i}$, the formula can be reduced to
\begin{equation}
  \begin{aligned}
    M^{C}_{r}(X_{i}; \lambda_{i}) &=
    \sum_{v \in V^{c}_{r}} \left(
    X_{i} \; - \; \left( X_{i} - \mathbf{d}_{i}(e(v))\right)^{+}
    \right) \\
    &= \sum_{v \in V^{c}_{r}} \left(
    X_{i} -
    \begin{cases}
      X_{i} - \mathbf{d}_{i}(e(v)) &\text{ if } X_{i} \ge \mathbf{d}_{i}(e(v)) \\
      0 &\text{ if } X_{i} < \mathbf{d}_{i}(e(v))
    \end{cases}
    \right) \\
    &= \sum_{v \in V^{c}_{r}} \left(
    \begin{cases}
      \mathbf{d}_{i}(e(v)) &\text{ if } X_{i} \ge \mathbf{d}_{i}(e(v)) \\
      X_{i} &\text{ if } X_{i} < \mathbf{d}_{i}(e(v))
    \end{cases}
    \right) \\
    &= \left|\{v: \mathbf{d}_{i}(e(v)) > X_{i}\}\right| \cdot X_{i} +
    \left|\{v: \mathbf{d}_{i}(e(v)) \le X_{i}\}\right| \\
    &= a_{0}(X_{i}) \cdot X_{i} + a_{1}(X_{i})
  \end{aligned}
  \label{eq:ck94-single-edge-memory-cost}
\end{equation}

\subsubsection{General Case: Multi-edge References}
The idea is to reuse \cref{eq:ck94-single-edge-memory-cost}. The distance vector
$\mathbf{d}(e(v))$ of only-edge $e(v)$ of reference $v$ is replaces with a
representative summarised distance $\hat{\mathbf{d}}(v)$.
\begin{equation}
  \hat{\mathbf{d}}_{i}(v) \coloneqq \min_{e \in E_{\text{in}(v)}} \mathbf{d}_{i}(e)
  \label{eq:ck94-multiedge-summary-distance}
\end{equation}

Note that in \cref{eq:ck94-multiedge-summary-distance}, each component of
$\hat{\mathbf{d}}(v)$ may come from different edges.

The summarised vector $\hat{\mathbf{d}}(v)$ represents the earliest possible
point in the unrolling space when a reference may be removed, so we could
choose the $X_{i}$ that is less than the actual optimum. So this is a
conservative edtimation.

\subsection{Invariant References}
Given an invariant reference $v \in V^{I}_{r}$ w.r.t. loop $\lambda_{k}$ ($1 \le
k < N$), where $N$ is the deepest level of loop nest (hence $\lambda_{N}$
denotes the innermost loop), we discuss the what we can gain by
unrolling a loop $\lambda_{k}$.

Recall that when a reference $v$ is invariant w.r.t. loop $\lambda_{i}$, all of
its array subscript expressions does not contain the induction variable of
$\lambda_{i}$.

In general, if we unroll $\lambda_{i}$, since $Invar(v, \lambda_{i})$
(the induction variable of $\lambda_i$ is not present in any array subscript
expression on $v$), the unrolled copies
of the same reference will access the same memory address (see
\cref{fig:ck94-outer-invariant-replacement}).
In one pass of the innermost loop, there is only one memeory access needed.

\begin{figure}[H]
  \centering
  \begin{minipage}[t]{\textwidth}
\begin{lstlisting}[language=C, caption={Original Loop}]
for (int i = 0; i < M; i++) {
  for (int j = 0; j < N; j++) {
    for (int k = 0; k < N; k++) {
      x = A[i][k]; /* R0 */
    }
  }
}
\end{lstlisting}
  \end{minipage}
  \par\medskip
  \begin{minipage}[t]{0.48\textwidth}
\begin{lstlisting}[language=C, caption={After Unroll-and-Jam}]
for (int i = 0; i < M; i++) {
  for (int j = 0; j < N; j += 3) {
    for (int k = 0; k < N; k++) {
      x = A[i][k]; /* R0 */
      y = A[i][k]; /* R1 */
      z = A[i][k]; /* R2 */
    }
  }
}
\end{lstlisting}
  \end{minipage}%
  \hfill
  \begin{minipage}[t]{0.48\textwidth}
\begin{lstlisting}[language=C,
  caption={Outer-loop invariant read after scalar replacement},
  label={lst:ck94-outer-invariant-replaced}]
for (int i = 0; i < M; i++) {
  for (int j = 0; j < N; j += 3) {
    for (int k = 0; k < N; k++) {
      tmp = A[i][k]; /* R0 */
      x = tmp;
      y = tmp;
      z = tmp;
    }
  }
}
\end{lstlisting}
  \end{minipage}
  \caption{Replacing repeated reads after unrolling an outer loop.}
  \label{fig:ck94-outer-invariant-replacement}
\end{figure}

If a reference is invariant with respect to the innermost loop, its load
can be hoisted outside that loop, as in
\cref{fig:ck94-inner-invariant-replacement}. In one pass of the innermost loop,
there is no memory read.

\begin{figure}[H]
  \centering
  \begin{minipage}[t]{\textwidth}
\begin{lstlisting}[language=C, caption={Original Loop}]
for (int i = 0; i < M; i++) {
  for (int j = 0; j < N; j++) {
    for (int k = 0; k < N; k++) {
      x = A[i][j]; /* R0 */
    }
  }
}
\end{lstlisting}
  \end{minipage}
  \par\medskip
  \begin{minipage}[t]{0.48\textwidth}
\begin{lstlisting}[language=C, caption={After Unroll-and-Jam}]
for (int i = 0; i < M; i++) {
  for (int j = 0; j < N; j += 3) {
    for (int k = 0; k < N; k++) {
      x = A[i][j]; /* R0 */
      y = A[i][j+1]; /* R1 */
      z = A[i][j+2]; /* R2 */
    }
  }
}
\end{lstlisting}
  \end{minipage}%
  \hfill
  \begin{minipage}[t]{0.48\textwidth}
\begin{lstlisting}[language=C,
  caption={Innermost-loop invariant reads after scalar replacement},
  label={lst:ck94-inner-invariant-replaced}]
for (int i = 0; i < M; i++) {
  for (int j = 0; j < N; j += 3) {
    tmp0 = A[i][j]; /* R0 */
    tmp1 = A[i][j+1]; /* R0 */
    tmp2 = A[i][j+2]; /* R0 */
    for (int k = 0; k < N; k++) {
      x = tmp0;
      y = tmp1;
      z = tmp2;
    }
  }
}
\end{lstlisting}
  \end{minipage}
  \caption{Hoisting a loop-invariant load outside the innermost loop.}
  \label{fig:ck94-inner-invariant-replacement}
\end{figure}

\begin{equation}
  \omega(v, i) =
  \begin{cases}
    0 &\text{if } Invar(v, \lambda_{N}) \\
    1 &\text{if } Invar(v, \lambda_{i}) \footnotemark \\
    X_{i} &\text{ otherwise}
  \end{cases}
  % \label{eq:}
\end{equation}

\footnotetext{This is an else-if case where $\neg Invar(v, \lambda_{N})$
is assumed.}


To analyse unrolling a specific loop $\lambda_{i}$
\begin{equation}
  \begin{aligned}
    M^{I}_{r}(X_{i}; \lambda_{i}) &= \sum_{v \in V^{I}_{r}} \omega(v, i)
    \\
    &= \left|\{v: \neg Invar(v, \lambda_{N}) \land
      Invar(v, \lambda_{i})\}\right| \\
    &\quad + \left|\{v: \neg Invar(v, \lambda_{i}) \land
      \neg Invar(v, \lambda_{N})\}\right| X_{i} \\
    &= a_{0} X_{i} + a_{1}
  \end{aligned}
\end{equation}

Note that here, $a_{0}$ and $a_{1}$ is dependent on $\lambda_{i}$, but not
directly on the value of $X_{i}$.

% For each reference, $v$, and the unrol factor of the loop nest, the number of 
The total number \todo{finish}
\begin{equation}
  M^{I}_{r}(\mathbf{X}) = \sum_{v \in V^{I}_{r}} \left(
    \prod_{1 \le i < N} \omega(v, i)
\right)
\end{equation}

% For each reference $v$, its memory access count is calculated as \todo{ref}
% \begin{equation}
% \prod_{1 \le i < N} \omega(v, i)
% \end{equation}
%
% Notice that as long as the reference is invariant w.r.t. the innermost loop, we
% can always find a way to reduce its memory access by hoisting the 

\subsection{Loop Balance Properties}
When we only consider unrolling loop $\lambda_{i}$, we replace the $\prod_{1 \le i
< N} X_{i}$ term with just $X_{i}$, so the loop balance $\beta_{L}(X_{i}; \lambda_{i})$
can be calculated as
\begin{equation}
  \begin{aligned}
    \beta_{L}(X_{i}; \lambda_{i})
    &= \frac{M^{\varnothing}(X_i;\lambda_{i}) + M^{C}_{r}(X_i;\lambda_{i}) +
      M^{I}_{r}(X_i;\lambda_{i})}{f \cdot X_{i}} \\
    &= \frac{
      a_{0} X_{i} + (a_{1}(X_{i}) \cdot X_{i} + a_{2}(X_{i})) +
      (a_{3} X_{i} + a_{4})
    }{f \cdot X_{i}} \\
    &= \frac{c_{0}(X_{i}) \cdot X_{i} + c_{1}(X_{i})}{f \cdot X_{i}}
  \end{aligned}
\end{equation}

Note that $\forall X_{i}: c_{0}(X_{i}), c_{1}(X_{i}) \ge 0$

Also we have the properties
\begin{itemize}[nosep]
  \item $\beta_{L}$ is non-increasing in the $i$-th dimension, and
  \item untilling multiple loop nests in an outermost- to innermost-loop
    order does not increase $\beta_{i}$
\end{itemize}

\section{Register Pressure}
The registers are used to (1) store intermediate subcsript calculation and (2)
hold values of regerences removed by scalar replacement.


\subsection{Non-dependent References}
For non-dependent references, no scalar replacement is made, so we only need to
model registers used in subscript expression.

To compute the number of registers used to store intermediate results for
subscript expression evaluation, the tree-labeling technique
\cite{sethiUllman1970} is used. However, it does not take into account the
increase in register pressure due to instruction scheduling
\cite{carrKennedy1994}.

\subsection{Variant References}
\Cref{eq:ck94-removed-variant-references} computes the number of
references removed by scalar replacement after unroll-and-jam.

\subsubsection{Spetial Case: Single-edge References}
We assuming all register interfere, all values live throughout the (inner?) loop body, and registers are not 
to be reused.

For each edge $e$, the number of registers
needed after scalar replacement is
\begin{equation}
\mathbf{d}_{N}(e) + 1
\end{equation}
where $\mathbf{d}_{N}$ is the component corresponding to the innermost loop
$\lambda_{N}$.

\todo{Read \cite{sethiUllman1970}}

For example, in \cref{lst:ck94-register-variant-original}, we have
$\mathbf{d}_{W_0\rightarrow V_0} = \left\langle 1, 2\right\rangle$, when
unrolling $\lambda_1$ (induction variable $i$) by $X_1 = 2$, we will have $\mathbf{d}_{W_0\rightarrow
V_1} = \left\langle 0, 2 \right\rangle$, which is \textit{carried by the
innermost loop}.
We could perform scalar replacement on this dependency edge to eliminate load.

In \cref{lst:ck94-register-variant-replaced}, two temporary \texttt{old} values
have been created to cache the values of $A[i][j-2], A[i][j-1]$, one extra
\texttt{current} value is created to cache the current $A[i][j]$ value.

So we need $\mathbf{d}_{N} + 1 = 2 + 1 = 3$ temporary values (registers).

\begin{figure}[H]
  \centering
  \begin{minipage}[t]{\textwidth}
\begin{lstlisting}[language=C,
  caption={Original loop with a variant reference},
  label={lst:ck94-register-variant-original}]
for (int i = 1; i < M; i++) {
  for (int j = 2; j < N; j++) {
    A[i][j] /* W */ = i + j;
    B[i][j] = A[i-1][j-2] /* V */ + 1;
  }
}
\end{lstlisting}
  \end{minipage}
  \par\medskip
  \begin{minipage}[t]{0.48\textwidth}
\begin{lstlisting}[language=C, basicstyle=\ttfamily\scriptsize,
  caption={Variant-reference loop after unrolling},
  label={lst:ck94-register-variant-unrolled}]
for (int i = 1; i < M; i += 2) {
  for (int j = 2; j < N; j++) {
    /* Original */
    A[i][j] /* W0 */ = i + j;
    B[i][j] = A[i-1][j-2] /* V0 */ + 1;
    /* Unrolled */
    A[i+1][j] /* W1 */ = i + 1 + j;
    B[i+1][j] = A[i][j-2] /* V1 */ + 1;
  }
}
\end{lstlisting}
  \end{minipage}%
  \hfill
  \begin{minipage}[t]{0.48\textwidth}
\begin{lstlisting}[language=C, basicstyle=\ttfamily\scriptsize,
  caption={Unrolled loop after scalar replacement},
  label={lst:ck94-register-variant-replaced}]
for (int i = 1; i < M; i += 2) {
  double old1 = A[i][0];
  double old2 = A[i][1];
  for (int j = 2; j < N; j++) {
    double current = i + j;
    B[i][j] = A[i-1][j-2] + 1;
    A[i+1][j] = i + 1 + j;
    /* eliminated load */
    B[i+1][j] = old1 + 1;
    /* shift values */
    old1 = old2;
    old2 = current;
  }
}
\end{lstlisting}
  \end{minipage}
  \caption{Unrolling and scalar replacement for a variant reference.}
  \label{fig:ck94-variant-register-example}
\end{figure}


\begin{equation}
  R^{C}_{r}(\mathbf{X}) =
  \sum_{v \in V^{C}_{r}} \left(
    \left(\prod_{1 \le i < N}
      \left(X_{i} - \mathbf{d}_{i}(e_{v})\right)^{+}\right)
    \cdot \left(\mathbf{d}_{N}(e_{v}) + 1\right)
    \right)
\end{equation}

\subsubsection{General Case: Multi-edge References}


\subsection{Invariant References}
For an invariant reference, \textit{repeated uses can share the one retained 
value}. Hence the register count depends on \textit{how many distinct values the
unrolled body needs}, not how many iterations that seperates the uses as opposed
to variant references.

\subsubsection{Spetial Case: Single-edge References}
When unroll loop $\lambda_i$ by factor $X_i$, we consider the following cacse.

\paragraph{Case $\neg Invar(v, \lambda_i)$} When the reference is variant with
the unrolling loop, \todo{}
\begin{figure}[H]
  \centering
  \begin{minipage}[t]{0.48\textwidth}
\begin{lstlisting}[language=C, basicstyle=\ttfamily\scriptsize,
  aboveskip=2pt, belowskip=2pt,
  caption={Repeated reads of A with $j$ advanced by two},
  label={lst:ck94-reg-invar-variant-j-step-two}]
for (int i = 0; i < M; i++) {
  for (int j = 0; j < N; j += 2) {
    use(A[i]);
    use(A[i]);
  }
}
\end{lstlisting}
  \end{minipage}%
  \hfill
  \begin{minipage}[t]{0.48\textwidth}
\begin{lstlisting}[language=C, basicstyle=\ttfamily\scriptsize,
  aboveskip=2pt, belowskip=2pt,
  caption={Cached reads of A after unrolling $i$},
  label={lst:ck94-reg-invar-variant-i-unrolled-cached}]
for (int i = 0; i < M; i += X1) {
  t_0 = A[i]; /* ... */ t_X1 = A[i];
  for (int j = 0; j < N; j += 2) {
    use(t_0); /* ... */ use(t_X1);
    use(t_0); /* ... */ use(t_X1);
  }
}
\end{lstlisting}
  \end{minipage}
  \caption{Caching repeated $A[i]$ reads after unrolling the outer loop.}
  \label{fig:ck94-reg-invar-variant}
\end{figure}

\paragraph{Case $Invar(v, \lambda_i) \land 
\exists j < N: \left(Invar(v, \lambda_j) \land X_{j} > 1\right)$} Hello

\begin{figure}[H]
  \centering
  \begin{minipage}[t]{0.48\textwidth}
\begin{lstlisting}[language=C, basicstyle=\ttfamily\scriptsize,
  aboveskip=2pt, belowskip=2pt,
  caption={Reads of A with $j$ advanced by two},
  label={lst:ck94-reg-invar-outer-j-step-two}]
for (int i = 0; i < N1; i++) {
  for (int j = 0; j < N2; j += 2) {
    for (int k = 0; k < N3; k++) {
      use(A[k]); use(A[k]);
    }
  }
}
\end{lstlisting}
  \end{minipage}%
  \hfill
  \begin{minipage}[t]{0.48\textwidth}
\begin{lstlisting}[language=C, basicstyle=\ttfamily\scriptsize,
  aboveskip=2pt, belowskip=2pt,
  caption={Shared value from A after unrolling $i$},
  label={lst:ck94-reg-invar-outer-i-unrolled-shared}]
for (int i = 0; i < N1; i += X1) { 
  for (int j = 0; j < N2; j += 2) {
    for (int k = 0; k < N3; k++) {
      t = A[k];
      use(t); use(t);
      /* ... */
      use(t); use(t);
    }
  }
}
\end{lstlisting}
  \end{minipage}
  \caption{Sharing an $A[k]$ read across copies of the outer loop.}
  \label{fig:ck94-reg-invar-outer}
\end{figure}

\paragraph{Case $Invar(v, \lambda_i) \land Invar(v, \lambda_N)$} Hello
\begin{figure}[H]
  \centering
  \begin{minipage}[t]{0.48\textwidth}
\begin{lstlisting}[language=C, basicstyle=\ttfamily\scriptsize,
  aboveskip=2pt, belowskip=2pt,
  caption={Reads of A with $k$ advanced by two},
  label={lst:ck94-reg-invar-inner-k-step-two}]
for (int i = 0; i < N1; i++) { 
  for (int j = 0; j < N2; j++) {
    for (int k = 0; k < N3; k += 2) {
      use(A[j]); use(A[j]);
    }
  }
}
\end{lstlisting}
  \end{minipage}%
  \hfill
  \begin{minipage}[t]{0.48\textwidth}
\begin{lstlisting}[language=C, basicstyle=\ttfamily\scriptsize,
  aboveskip=2pt, belowskip=2pt,
  caption={Load from A hoisted after unrolling $i$},
  label={lst:ck94-reg-invar-inner-i-unrolled-hoisted}]
for (int i = 0; i < N1; i += X1) {
  for (int j = 0; j < N2; j++) {
    t = A[j];
    for (int k = 0; k < N3; k++) {
      use(t); use(t);
      /* ... */
      use(t); use(t);
    }
  }
}
\end{lstlisting}
  \end{minipage}
  \caption{Hoisting an $A[j]$ read out of the innermost loop.}
  \label{fig:ck94-reg-invar-inner}
\end{figure}

\paragraph{Other cases} Not considered for scalar replacement
% \begin{figure}[H]
%   \centering
%   \begin{minipage}[t]{0.72\textwidth}
% \begin{lstlisting}[language=C, basicstyle=\ttfamily\scriptsize,
%   aboveskip=2pt, belowskip=2pt,
%   caption={Read invariant only in the middle loop},
%   label={lst:ck94-reg-invar-middle-only}]
% for (int i = 0; i < N1; i++) { 
%   for (int j = 0; j < N2; j++) {
%     for (int k = 0; k < N3; k++) { /* <- invariant */
%       use(A[i][k]);
%     }
%   }
% }
% \end{lstlisting}
%   \end{minipage}
%   \caption{An invariant reference not covered by the preceding cases.}
%   \label{fig:ck94-reg-invar-other}
% \end{figure}

\paragraph{}
Let $\alpha(v, i; \mathbf{X})$ denote the nubmer of registers needed to do
scalar replacement after expanding loop $\lambda_i$ by $X_i$, given the
unrolling factor of other loops $\mathbf{X}$.
\begin{equation}
  \alpha(v, i; \mathbf{X}) = 
  \begin{cases}
    X_i &\text{if } \neg Invar(v, \lambda_i) \\
    1 &\text{if } \left(\exists j < N: Invar(v, \lambda_j) \land X_j > 1\right)
           \lor \left(Invar(v,\lambda_N)\right)
           \footnote{This is an else-if condition where $Invar(v, \lambda_i)$ is
           assumed.}
           \\
    0 &\text{Otherwise}
  \end{cases}
  \label{eq:}
\end{equation}

For a given reference, we can model the number of registers needed after
unroll-and-jam and scalar replacement by 
\cref{eq:reg-invariant-total-to-ref}
\begin{equation}
  \prod_{1 \le i \le N} \alpha(v, i; \mathbf{X}) 
  \label{eq:reg-invariant-total-to-ref}
\end{equation}

Notice in \cref{eq:reg-invariant-total-to-ref}, the "otherwise" case
produces 0, which results a total of 0 in the product. This means that in the
specific case of analysing registers added after scalar replacement for
invariant references, as long as $\mathbf{G}(v, \mathbf{X})$ is not satisfied.
\begin{equation}
  \begin{aligned}
    1 - \mathbf{G}(v \in V^{I}_{r}, \mathbf{X}) &\coloneqq
    \left(
      \begin{aligned}
        \exists 1 \le i \le N:\quad & Invar(v, \lambda_i) \\
          & {}\land \neg Invar(v, \lambda_N) \\
          & {}\land \left(
            \begin{aligned}
              \forall j < N:\quad & \neg Invar(v, \lambda_j) \\
                & {}\lor (Invar(v, \lambda_j) \land X_j = 1)
            \end{aligned}
          \right)
      \end{aligned}
    \right) \\
    &\equiv
    \left(
      \begin{aligned}
        &(\exists 1 \le i \le N:\, Invar(v, \lambda_i)) \\
        &{}\land \neg Invar(v, \lambda_N) \\
        &{}\land (\forall j < N:\, Invar(v, \lambda_j) \implies X_j = 1)
      \end{aligned}
    \right) \\
    &\overset{(*)}{\equiv}
    \left(
      \begin{aligned}
        &\neg Invar(v, \lambda_N) \\
        &{}\land \forall j < N:\, (Invar(v, \lambda_j) \implies X_j = 1)
      \end{aligned}
    \right)
  \end{aligned}
  \label{eq:reg-invar-ref-cond}
\end{equation}
Here, $(*)$ uses the assumption
$v \in V^{I}_{r} \implies \exists i:\, Invar(v, \lambda_i)$. 

\Cref{eq:reg-invar-ref-cond} shows that, if a loop is variant to the innermost
loop and for every of its invariant loops, we choose not to unroll, then there
is no opportunity for the scalar replacement to reduce any memory access (by
caching values in a register), as illustrated by
\cref{lst:ck94-no-scalar-replacement}.

\begin{lstlisting}[language=C,
  caption={Reference varying in the innermost loop},
  label={lst:ck94-no-scalar-replacement}]
for (int i = 0; i < M; i++) { /* <- no unrolling for its invariant loop */
  for (int j = 0; j < N; j++) { /* <- variant w.r.t. innermost loop */ 
    use(A[j]); /* introducing register does NOT reduce memory access */
  }
}
\end{lstlisting}

